% --- Άσκηση 34152 (34152.tex) ---
\Askhsh{34152}{2}
\begin{ylh}{2}
    \item 2.3. Απόλυτη Τιμή Πραγματικού Αριθμού
    \item 2.4. Ρίζες Πραγματικών Αριθμών
\end{ylh}
Δίνονται οι παραστάσεις: $A = \sqrt{(x - 2)^2}$ και $B = \sqrt[3]{(2 - x)^3}$, όπου $x$ πραγματικός αριθμός.

\begin{alist}
\item Για ποιες τιμές του $x$ ορίζεται η παράσταση $A$; \monades{7}
\item Για ποιες τιμές του $x$ ορίζεται η παράσταση $B$; \monades{8}
\item Να δείξετε ότι, για κάθε $x \le 2$, ισχύει $A = B$. \monades{10}
\end{alist}
